\documentclass[11pt]{article}
\usepackage{fullpage}
\usepackage{graphicx}
\usepackage{url}

\title{The Architecture of FindBugs}
\author{David Hovemeyer}
\date{}

\newcommand{\Note}[1]{
\begin{quote}
\fbox{
\parbox{.5in}{
\includegraphics[width=.4in]{attention}
}
\parbox{5in}{
{\bf Note}: #1
}
}
\end{quote}
}

\newcommand{\Indent}{\hskip .25in}
\newcommand{\ForEach}{{\bf for each}}
\newcommand{\Do}{{\bf do}}
\newcommand{\If}{{\bf if}}
\newcommand{\Then}{{\bf then}}
\newcommand{\EndFor}{{\bf end for}}
\newcommand{\EndIf}{{\bf end if}}

\begin{document}

\maketitle
\tableofcontents

\section*{Introduction}

This document presents a medium-level overview of the architecture
of FindBugs.  It should be helpful for anyone trying to understand how
the tool works and how to extend its capabilities.

This information is accurate as of FindBugs version 0.9.4.
The document revision is
\begin{quote}
\verb+$Revision: 1.4 $+
\end{quote}

Note that FindBugs has fairly complete javadoc API documentation.
The {\tt apiJavadoc} target in the Ant build.xml file will generate
this documentation in the {\tt apiJavaDoc} subdirectory.

\section{Walkthrough}

This section will provide a walkthrough of an analysis session,
such as an invocation of the command line interface, or the analysis
started by the ``Find Bugs!'' button in the Swing GUI.

\subsection{The FindBugs class}

%The FindBugs class serves as both the main engine and the command line
%user interface.  When the {\tt bin/findbugs} script is run, most of the
%command line arguments are passed along verbatim to the main() method of
%the FindBugs class.

First, an instance of the FindBugs class is created.
The FindBugs class serves as the main engine that sets up and executes the
analysis.  The FindBugs constructor takes
two parameters: a BugReporter and a Project.

The BugReporter receives all warnings generated by the analysis.
Several bug reporter implementations exist in order to support various
output formats.  The XMLBugReporter class is perhaps the most useful;
it can save generated warnings in a format that can be reloaded without
loss of information.  

A Project represents the software artifact to be analyzed.
The main information in the Project is a list of files and directories
containing classes to analyze.  These classes are the ``application classes''.
The Project also contains a list of
files and directories containing classes that are {\em referenced}
by the application classes, but which should not themselves be analyzed.
These are the ``auxiliary classpath'' classes.  The Project contains
a list of source directories and archives: this list can be consulted to
find the source code for a range of code reported in a warning.

\subsection{Configuring the FindBugs object}

There are lots of ways for the user to configure the analysis parameters.
Some examples:
\begin{itemize}
\item Which detectors are enabled and disabled
\item Default reporting threshold (low, medium, high)
\item Which bug categories to report
\item Rule-based filtering of warnings
\item The analysis effort level: greater effort, more bugs found and longer running time
\end{itemize}

These parameters are generally configured using accessor methods which
modify fields in the FindBugs object.  The UserPreferences class serves
to organize many of these parameters into a single object.

\subsection{FindBugs.execute()}

Once the FindBugs object has been created and configured, the execute() method
starts the analysis.

\subsubsection{The Analysis Context}

First, an AnalysisContext object is created.  This object is the repository for
all global information pertaining to the current analysis session.
Examples of information stored in the analysis context:
\begin{itemize}
\item The Subtypes object representing the class hierarchy
\item Databases of interprocedural analysis facts
\end{itemize}

In addition, a number of methods supporting access to the classpath and the
repository of application classes are available in the AnalysisContext class.
Although these tasks can be done using static methods in the BCEL
Repository class, eventually we would like to decouple FindBugs from BCEL
and its deeply flawed Repository implementation.  Therefore, any classpath
or class hierarchy queries should be done via the AnalysisContext.

Another way to look at AnalysisContext is that it allows FindBugs to do
whole-program analysis.  Most FindBugs analysis is local: but as deeper
analysis is implemented in FindBugs (e.g., interprocedural analysis,
whole-program call graph construction, etc.), this functionality should
be available from the AnalysisContext.

Originally, the reference to the AnalysisContext object was explicitlyJennifer Polack <polack@umw.edu>

passed to all bug detector objects using a setAnalysisContext() method.
More recently, the AnalysisContext is stored in an InheritableThreadLocal, and is accessible
from a static method, AnalysisContext.currentAnalysisContext().

Once the AnalysisContext has been created, the ``analysis features'' are configured.
Analysis features are configuration properties stored in the AnalysisContext that
affect how the analysis is performed.

\subsubsection{Analysis Features and Training Databases}

Around the same time as configuration of analysis features, training
databases are created and configured.  A training database records analysis facts
pertaining to a particular method or field, and allow information collected
in one analysis run to be re-used in a later run.  For example, methods that
unconditionally dereference method parameters can be recorded in one run,
and loaded in a later run.  Any time a null value is passed for an unconditionally
dereferenced parameter, FindBugs can issue a warning about a possible null pointer
dereference.  Training databases are thus a simple way to do limited interprocedural
analysis.

\subsubsection{Execution Plan}

The next main task of the execute() method is to create an execution plan.
FindBugs has recently acquired the ability to define ordering constraints
among bug detectors, including the ability to specify that a detector must
run in an earlier or later pass than another detector.
The execution plan explicitly organizes bug detectors into passes,
and orders the execution of detectors within passes.
Analysis passes can be used in much the same way as training databases:
information about methods and fields is collected in the first pass,
and then used in the second pass.  The advantage over training databases
is that both passes run in the context of a single
invocation of FindBugs.
It is also sometimes useful to specify that within a single analysis pass,
one detector runs before another.  With a single pass, each analyzed class
is considered (in no particular order).  The detectors in the pass are
applied to the class in order.  Therefore, by specifying an ordering
constraint between two detectors, the earlier detector can collect information
on the class that can then be used by a later detector (or detectors).

\Note{
The implementation of ExecutionPlan and related classes is needlessly complicated,
and should probably be rewritten at some point.
}

\subsubsection{Scanning for Application Classes}

After the execution plan has been created, the next task of execute() is
to determine which classes will be analyzed.  This is done by scanning
all class files, archive files, and directories specified to be part of
the analyzed application by the Project.  As each class is scanned,
it is parsed into a BCEL JavaClass object and placed in the BCEL Repository
(via the AnalysisContext).
Thus, all application classes are in memory during the entire analysis.
(This is one reason why FindBugs requires a lot of memory.)

\Note{
The BCEL Repository is implemented using static fields and static methods.
This makes FindBugs non-reentrant, so it is not possible to create two
FindBugs objects and perform two separate analysis sessions at the same time.
Unfortunately, BCEL uses the static Repository extensively in its
internal implementation, meaning that the Repository is very tightly coupled
with the rest of of BCEL.
The only long-term solution to this problem is to drop BCEL.
(Which we should do anyway.)
}

One subtlety of the class scanning process is that we look for things like
Class-Path attributes in Jar manifests in order to discover referenced
libraries.  When analyzing Java bytecode, it is often necessary to inspect
classes that are referenced by analyzed code in order to make the analysis
as accurate as possible.  For example, when analyzing a method call,
you can't tell what exceptions can be thrown unless you can examine the
classfile of the class specified as the receiver of the method.  Another reason
that referenced classes are necessary is that any analysis of the inheritance
hierarchy (including the analysis needed to perform type inference on
class files) requires access to referenced classes.  Thus, FindBugs wants
the auxiliary classpath to be as complete as possible.  The additional codebases
discovered from class scanning are themselves examined for additional classpath
references: in the end, we should have the transitive closure of codebases
referenced in this way.  Hopefully, the user has also made an effort to
explicitly specify referenced classes.

\subsubsection{Performing the Analysis}

\begin{figure}
\begin{tabbing}
\ForEach\  analysis pass in the execution plan \Do \\
\Indent \= \ForEach\  application class \Do \\
        \> \Indent \= \ForEach\  detector in the analysis pass \Do \\
        \>         \> \Indent \= create a ClassContext for the class\\
        \>         \>         \> apply the detector to the ClassContext\\
        \>         \> \EndFor \\
        \>  \EndFor \\
\EndFor
\end{tabbing}
\caption{Performing the analysis using the execution plan.}\label{fig:performingTheAnalysis}
\end{figure}

Once the execution plan is in place and all of the application classes have been collected,
we're ready to perform the analysis.  This is done more or less according to the
pseudo-code shown in Figure~\ref{fig:performingTheAnalysis}.

Each bug detector is an object implementing the Detector interface,
which specifies a visitClassContext() method.  This method takes
an instance of ClassContext created from the analyzed class.  The ClassContext serves
as a cache for analysis results from a single class, because many
kinds of analysis (such as type inference) are used by many detectors.
By performing an analysis only once and then sharing it among multiple
detectors, we decrease CPU time (at the cost of increasing memory use).
(ClassContext is described in more detail in Section~\ref{sec:bytecodeAnalysisPackage}.)

\subsubsection{Finishing Up}

After all analysis passes have executed, we call finish() and
reportQueuedErrors() on the BugReporter, and then clear the BCEL Repository
(since we no longer need to access classes contained therein).

\section{Bug Detectors}

This section presents a high-level overview of bug detectors and
how they work.  The bug detectors implemented in FindBugs can be divided into
two basic categories: visitor-based and CFG-based.  Ultimately, FindBugs
does not place any constraints on how bug detectors are implemented:
any analysis technique that works can and should be used.  However,
the visitor-based detectors tend to be used for simpler ``peephole''
analyses, and the CFG-based detectors tend to be used for more
complex analyses, especially ones that use the dataflow analysis
framework.

Ultimately, the task of a bug detector is very straightforward:
look at a single classfile and find potential bugs in it.
Potential bugs are reported by creating a warning object (instance
of the BugInstance class), decorating it with information about
the context of the warning and the reasons for reporting it,
and then reporting it via the BugReporter object.

All of the ``core'' bug detectors---those that we include with FindBugs
and are enabled by default---live in the edu.umd.cs.findbugs.detect
package.  However, FindBugs has a plugin architecture that allows users
to create their own detectors and use them without modifying FindBugs.
By dropping a FindBugs plugin (which is just a jar file with an XML descriptor,
message translation files, and detector classes) in the
{\tt plugin} directory of a FindBugs installation, new detectors can
easily be added.  For example, FindBugs core team member Dave Brosius maintains
the fb-contrib project (\url{fb-contrib.sourceforge.net}), which
provides additional bug detectors.

\subsection{Visitor-Based Detectors}

The visitor-based detectors are generally subclasses of BytecodeScanningVisitor,
which itself is a subclass of DismantleBytecode.  The base classes perform
a top-down traversal of the features of the classfile, decoding all
of the symbolic information.  When a feature of the classfile is encountered,
such as a field, method, instruction, and so forth, a callback method is
invoked.  By overriding these callback methods, a visitor-based detector
can inspect the classfile looking for suspicious features that may
indicate bugs.  For example, the FindUnconditionalWait detector looks for
methods containing a {\tt monitorenter} instruction followed by an {\tt invokevirtual}
instruction that invokes {\tt Object.wait()}.

One important idiom used in visitor-based detectors is a state machine
recognizer over the sequence of instructions in a method.
Each invocation of {\tt sawOpcode()}, which is the callback method
that handles individual instructions, is a single input symbol to
the state machine.  The overall state machine is essentially a finite-state
automaton that accepts some regular language, where the language is a
pattern that, if it appears in the bytecode for a method, indicates
a probable bug.  In other words, if the string of input symbols (instructions)
causes the state machine to accept the string, then a warning is issued.
Even though this kind of analysis is very simple,
especially in its simplistic treatment of control flow, detectors implemented
this way can be surprisingly effective for many kinds of bugs.

In addition to being simple and concise, visitor-based detectors
tend to be {\em much} faster than the CFG-based detectors.  Back when FindBugs
used visitor-based detectors exclusively, it was possible to analyze a large
jar file such as {\tt rt.jar} in several seconds.

A recent improvement to the visitor-based detectors is the OpcodeStack
class, which attempts to maintain information about the operand stack as
each instruction in a method is visited.  However, it does this in a fairly
ad-hoc way.  In the future we may implement something more sophisticated
to take the place of OpcodeStack.  (For example, we could use the existing
dataflow analyses to populate the operand stack contents, providing a bridge
from the visitor-based detectors to the more sophisticated CFG-based
analyses.)

\subsection{CFG-Based Detectors}

CFG stands for ``Control Flow Graph''.  The CFG-based detectors use a
CFG representation of Java methods to perform somewhat more sophisticated
analysis than the visitor-based detectors.  In particular, FindBugs
implements many kinds of dataflow analysis that can be used in CFG-based
bug detectors.

Unlike the visitor-based detectors, which generally inherit from
BytecodeScanningVisitor,
CFG-based detectors generally just implement the Detector interface
directly.

\begin{figure}
\begin{tabbing}
\ForEach\ method in the class \Do\\
\Indent \= request a CFG for the method from the ClassContext \\
        \> request one or more analysis objects on the method from the ClassContext \\
        \> \ForEach\ location in the method \Do \\
        \> \Indent \= get the the dataflow facts at the location \\
        \>         \> inspect the dataflow facts \\
        \>         \> \If\ a dataflow fact indicates an error \Then \\
        \>         \> \Indent \= report a warning \\
        \>         \> \EndIf \\
        \> \EndFor \\
\EndFor
\end{tabbing}
\caption{Pseudo-code for a CFG-based detector.}\label{fig:cfgBasedDetector}
\end{figure}

Figure~\ref{fig:cfgBasedDetector} shows the pseudo-code for the
visitClassContext() method of a CFG-based detector.
The basic idea is to visit each method of an analyzed class in turn,
requesting some number of analysis objects.  An analysis object
is simply the end product of some analysis: it records facts 
(or probable facts) about the method based on some form of program
analysis, usually dataflow analysis.
After getting the required analyses, the detector iterates through
each location in the control flow graph.  A location is
the point in execution just before a particular instruction is executed
(or after the instruction, for backwards analyses).  Locations are not quite synonymous
with instructions because of {\tt jsr} subroutines---more on those
in Section~\ref{sec:bytecodeAnalysisPackage}.  At each location,
the detector checks the dataflow facts to see if anything suspicious is
going on.  For example, the null dereference detector uses a null value
dataflow analysis to see if any instructions that have a runtime null
check are used on a value that is either definitely null, or is null
on a path that through the CFG that is likely to be feasible.
If suspicious facts are detected at a location, then the detector issues a warning.

CFG-based detectors are useful for any kind of analysis that can be
defined as a forward or backward dataflow problem.  Unfortunately,
they tend to be far slower than the visitor-based detectors, but hey,
computers are fast.\footnote{I am not sure if we know exactly why the
CFG-based detectors are so much slower.  They do require BCEL to create
heavier-weight representations of the class, methods, and instructions,
which takes up time and memory.  Constructing the CFG and performing the
dataflow analyses takes additional time.  I suspect that at least a 2x or 3x
speedup is possible.}

\section{The Bytecode Analysis (BA) Package}\label{sec:bytecodeAnalysisPackage}

The edu.umd.cs.findbugs.ba package is the ``bytecode analysis'' package
(hereafter referred to as the ``BA package''),
and is intended to be a generic library of classes for doing static analysis
of Java classes.  As such, the classes defined in this package (and subpackages)
should not refer to any of the ``bug-finding'' classes, such as the
FindBugs warning object model or detector classes.  References to BCEL
and the FindBugs generic graph library are necessary and permitted.

\Note{
Although parts of the bytecode analysis package, such as the CFG builder,
might be useful in other contexts, most of the analyses implemented in
the BA package are probably not that useful in other contexts.
In particular, no analysis used in FindBugs is designed to be {\em sound},
meaning that it deduces only facts that are true.  Many, many liberties are
taken by the analyses implemented in the BA package: while useful for
finding probable bugs, you wouldn't want to use these analyses in a compiler
or in a program verifier for safety critical systems.
}

\Note{
Until recently the BA package was a single flat Java package.  Due to
the large number of classes, we have started to organize it into subpackages.
This work is still ongoing.  Also, you may find that there are some unused
or experimental subpackages, such as the ``type2'' subpackage.
}

\subsection{AnalysisContext and ClassContext}

The AnalysisContext and ClassContext classes described earlier are part
of the BA package.  An AnalysisContext object is first created,
and its repository populated with classes to analyze (as well as
the codebases of referenced classes).  Analyses may be performed
on a particular class by submitting the JavaClass representing the
class to the AnalysisContext, getting back the ClassContext for the
class, and then requesting particular analyses directly from the
ClassContext.

Initially, a ClassContext has only a reference to the JavaClass object.
Particular analyses are performed on demand and then cached.  A few of
the analyses, such as the AssignedFieldMap, are global to the class
as a whole, and are obtained via a no-argument method.  Most of the
analyses, however, are performed on methods individually, and are obtained
via methods that take a BCEL Method object as a parameter.  In any case,
whenever an analysis is requested from the ClassContext a cache is first checked.
If the analysis has already been performed, the cached version is returned.
Otherwise, the analysis is performed, added to the cache, and returned.

Almost all methods that request an analysis can throw both CFGBuilderException
and DataflowAnalysisException.  Therefore, detectors using the BA package will generally need to
be able to handle these exceptions.  A typical strategy is to surround the
handling of each method in a detector with a try/catch that catches these exceptions and
reports them to the BugReporter if they occur.  Even if an exception occurs
in the checking of a single method, the detector should generally continue
with the other methods in the analyzed class.

\subsection{Control Flow Graphs}

The Control Flow Graph implementation consists of the CFG, BasicBlock, Edge,
and BetterCFGBuilder2\footnote{This class was preceeded by the
SimpleCFGBuilder and BetterCFGBuilder classes.  Constructing a CFG from
Java bytecode is surprisingly difficult.  BetterCFGBuilder2 has proven
to be very robust, so further evolution will probably not be needed.} classes.

CFG, BasicBlock, and Edge are a very this layer on top of the BCEL MethodGen,
InstructionHandle, and Instruction classes.  FindBugs uses BCEL InstructionHandle
objects to represent single instructions.  Thus, if you ask a BasicBlock
for the instructions it contains, it will give you back InstructionHandle
objects.

BetterCFGBuilder2 takes a MethodGen and uses it to construct a CFG object.
This process is complicated by the fact that Java methods may use
the {\tt jsr} instruction, usually the result of a try/finally construct.
Code reachable from a {\tt jsr} instruction is known as a ``jsr subroutine''.
FindBugs handles jsr subroutines by simply inlining them into the control flow graph.
Because they are guaranteed not to be recursive, this process is guaranteed to
terminate for any valid Java method.

Because of jsr subroutines, an individual bytecode instruction (InstructionHandle) may appear
multiple times in the CFG, once for each expansion of the subroutine it appears in.
This is the reason for the Location class: a Location is a BasicBlock, InstructionHandle
pair.  All dataflow analyses use Location objects as a key to retrieve computed
analysis facts.  The CFG class has a locationIterator() method that returns an
Iterator over all defined locations in a CFG.

Another subtlety of FindBugs CFGs is how exception control flow is
represented.  Generally, any instruction that can throw an exception is
preceeded by an Exception Throwing Block (ETB).  The exception control
edge coming from an ETB represents what will happen if the exception occurs:
either transferring control to an exception handler, or throwing it out
of the method.  The fall-through edge coming from the ETB represents
what will happen if no exception occurs.  One surprising feature of
ETB blocks is that they appear to be empty: that is, they contain no
actual instruction.  Instead, they represent an implicit ``conditionally throw an
exception'' instruction: for example, a null check.  ETBs may be
identified by calling the isExceptionThrower() method on a BasicBlock.

Initially, a FindBugs CFG will assume that any ETB can transfer control
to any exception handler in the method.  This conservative assumption results
from the fact that, during CFG construction, we don't really know what
methods are being called, and what exceptions those methods can throw.
A CFG with conservative exception edges is a ``raw'' CFG---you can see
some evidence of raw CFGs in ClassContext, and a few other places.
To the extent possible, ClassContext tries to convert raw CFGs into 
``pruned'' CFGs by removing obviously infeasible exception edges.
The more accurate the exception edges, the more accurate analyses that
use CFGs can be.

\subsection{Dataflow Analysis}

FindBugs implements a very generic dataflow analysis framework.
It is generic in the sense that any kind of object can be used as
a dataflow fact, and also in the sense that type parameters are used
whereever appropriate.  Thus, adding a new kind of dataflow analysis
to FindBugs tends to be a fairly simple exercise.

This section will briefly discuss the generic parts of the dataflow framework,
and also discuss some of the most important concrete dataflow analysis
classes.

\subsubsection{Dataflow and DataflowAnalysis}

The DataflowAnalysis interface is the supertype for all concrete dataflow
analysis classes.  It defines methods for creating and manipulating dataflow
facts.  The Dataflow class implements the actual dataflow analysis algorithm
using a CFG and an instance of DataflowAnalysis.

AbstractDataflowAnalysis is an abstract class implementing DataflowAnalysis,
and providing implementations of some of its methods.  In addition,
AbstractDataflowAnalysis defines the notion of implementing the dataflow transfer
function at the level of individual instructions rather than complete basic
blocks, and automatically implements a cache of analysis facts at each
CFG location.  Most dataflow analyses will want to inherit from
AbstractDataflowAnalysis.  However, there are some exceptions.
For example, DominatorsAnalysis is a dataflow analysis where the transfer
function makes sense only at the level of complete basic blocks:
individual instructions have no effect on dataflow facts.

ForwardDataflowAnalysis and BackwardDataflowAnalysis, as their names suggest,
are generic base classes for forward and backward dataflow analyses.
Both of these classes inherit from AbstractDataflowAnalysis.
So, when defining a new dataflow analysis class, inheriting from one of these
classes is a reasonable choice.  The StackDepthAnalysis class is a very good
example of how to implement a forward dataflow analysis in FindBugs.
It uses the StackDepth class to represent dataflow facts, and computes the
depth of the Java operand stack at each location in the analyzed CFG.

\subsubsection{Frame, FrameDataflowAnalysis, and AbstractFrameModelingVisitor}

In many dataflow analyses, we want to know something about the
values of local variables and stack operands.  The Frame class represents
the Java stack frame at a single CFG location, and represents both
stack operands and local variables.  Note that both stack operands
and locals are considered to be ``slots'': there is really nothing
fundamentally different about them, other than the way that they
are accessed by instructions.  Most Java instructions operate exclusively
on stack operands.

\Note{
A peculiarity of Frame-based dataflow analysis in FindBugs is that
long and double values are considered to occupy two stack slots, or
two local variable slots.  This quirk is inherited from BCEL,
which inherits it from the original JVM specification from Sun.
}

The Frame class is parameterized
with the type of the values stored in the locals and stack locations;
generally, you should inherit from Frame to create a Frame subtype that
is specific to your analysis.  For example, IsNullValueFrame is a frame
subtype where the locals and stack operands are of type IsNullValue,
and is used (no surprise) by IsNullValueAnalysis.

FrameDataflowAnalysis is a convenient abstract base class for dataflow analysis
classes that will use Frames as dataflow facts.  Its specific purpose is
to allow the dataflow merge function to be specified on individual values
rather than entire frames.

AbstractFrameModelingVisitor is a useful base class for defining a visitor
that models the effect of bytecode instructions on a Frame of values.
Using the notion of a ``default'' dataflow value, it provides generic implementations
for all Java bytecode instructions.  You can then simply override the
instruction-visiting methods that are actually important for the particular
analysis.

\subsubsection{Using the Results of a Dataflow Analysis}

The most convenient way to access and use the results of applying
a particular dataflow analysis to a method is to get a reference to the
Dataflow object that was used to perform an analysis.  For example,
to get the results of applying IsNullValueAnalysis, get a reference to
the IsNullValueDataflow object that executed the analysis.  (You can get
this object directly from the ClassContext using the
getIsNullValueDataflow, passing the Method object you want to analyze.)

All Dataflow subclasses support the methods getStartFact() and getResultFact(),
which get the dataflow fact before or after a particular basic block, respectively.
Most Dataflow subclasses also support methods called getFactAtLocation() and
getFactAfterLocation(), which return the dataflow fact immediately before
or after a particular Location in the CFG.

\subsubsection{TypeAnalysis and TypeDataflow}

One of the most important dataflow analyses in FindBugs is TypeAnalysis/TypeDataflow.
This frame-based analysis performs type inference on a Java method, computing
types for all local variables and stack operands.

The dataflow merge in this analysis computes the least upper bound
of the merged values in the {\em class} inheritance hierarchy.
Note that the least upper bound for superinterfaces of two types
is not guaranteed to yield a unique result in Java, due to multiple inheritance
of interfaces.  Therefore, TypeAnalysis in FindBugs uses the same recourse
as the JVM verifier, which is to consider only superclasses.
One consequence of this limitation is that the type of a receiver object in an
{\tt invokeinterface} instruction is often only known as java.lang.Object,
because dataflow merges tend to lose information about implemented interfaces.

\Note{
In the long term, it would be nice to have type information that was based on
{\em sets} of types.
}

\subsubsection{ValueNumberAnalysis and ValueNumberDataflow}

A common problem in analysis of Java bytecode is determining when
two locations in a stack frame have the same value.  This problem is
especially important when you consider that Java bytecode is stack-based:
most instructions can only manipulate stack values, not values in
local variables.  Information gained when modeling an instruction
that affects values on the stack should also be applied to identical
values elsewhere on the stack, or in local variables.

ValueNumberAnalysis and ValueNumberDataflow attempt to assign value numbers
to each location in each frame.  If two frame slots have the same value number,
then at runtime then definitely contain the same value.

One application of value numbers is in the null dereference detector,
which makes use of the various {\tt if} comparison instructions:
when a value is compared against null, then we learn something about
the value: depending on which branch is taken, the value is either definitely
null or definitely not null.  The {\tt if} comparison instructions---{\tt ifnull},
{\tt ifnonnull}, {\tt ifacmp\_eq}, and {\tt ifacmp\_ne}---operate on
values that are on top of the operand stack.  On the edges leading from
these instructions, we can use the information gained about the value
compared to null on the top of the operand stack to update other occurrences of
the same value elsewhere on the operand stack, or in local variables.

One feature of the value number analysis in FindBugs is that it
performs redundant load elimination and forward substitution.
This is required in blocks of code like the following:
\begin{quote}
\begin{verbatim}
synchronized (this.lock) {
    while (!someCondition()) {
        this.lock.wait();
    }
}
\end{verbatim}
\end{quote}
%
Without redundant load elimination, the value number analysis would not
assume that the two references to {\tt this.lock} refer to the same value.
While the Java memory model does allow the two loads to return
different values (e.g., another thread might update the field between
accesses), by far the most likely behavior is that both loads will
return the same value (and most JVMs will optimize the code to
guarantee this outcome).  Without redundant load elimination, FindBugs
might assume that the {\tt wait()} call is made to a different object
than the one synchronized on.

\subsubsection{IsNullValueAnalysis and IsNullValueDataflow}

IsNullValueAnalysis determines which frame slots contain definitely-null
values, definitely non-null values, and various kinds of conditionally-null
or uncertain values.  This information is used by the FindNullDeref
detector to detect possible null pointer dereferences and redundant null
comparisons.

\subsubsection{LockAnalysis and LockDataflow}

LockAnalysis computes LockSet values at each CFG location.  A LockSet is
the set of value numbers that are known to be locked.  This analysis
is used, for example, by the FindInconsistentSync2 detector, which looks for
unsynchronized field accesses that might be reachable from multiple threads.
Note that this analysis is strictly intra-procedural: it does not consider
that some values might be locked in a calling method.

\subsection{Class Hierarchy Queries and Missing Classes/Interfaces}

A very common task in bug detectors and static analysis in general
is to perform class hierarchy queries.  Some examples:
\begin{enumerate}
\item Find immediate superclass and superinterfaces of a type
\item Find transitive superclasses and superinterfaces of a type
\item Find the superclass method that a particular method overrides or implements
\item Find all subtypes of a particular class or interface
\item Find all possible methods that might be called by a particular
      {\tt invokevirtual} instruction
\end{enumerate}

The first three examples are {\em upwards} queries: they only require
superclasses and superinterfaces to answer.  Since superclasses and superinterfaces
must be present for a class to be loaded by the JVM, it is reasonable
to assume that a static analysis can find them.  Therefore, we
make the assumption that upwards queries can always be answered definitively,
and that if a superclass or superinterface cannot be found, it is
a serious error.  In fact, many methods in BCEL implicitly do
upwards queries using the BCEL Repository class, and throw
ClassNotFoundException if a required superclass or superinterface
cannot be found.  Any time you are in a context where you need to
handle ClassNotFoundException, you should find a RepositoryLookupFailureCallback
object and pass the exception to its reportMissingClass() method.
All BugReporter objects implement RepositoryLookupFailureCallback,
so you can use the BugReporter.  You can also call the static
reportMissingClass() method in AnalysisContext, which will also locate
the current lookup failure callback (generally the bug reporter)
and pass the exception along.

The last two examples are {\em downwards} queries: they attempt to answer
questions about possible subtypes.  This is closed world assumption that
is generally unsound, because Java allows classes to be loaded dynamically.
However, in a static analysis to find bugs we don't require soundness,
and for most real applications it a reasonable assumption that the user
will give the analysis access to the complete application.  Therefore,
FindBugs uses both upwards and downwards class hierarchy queries.
(Upwards queries are used much more extensively.)

It would be nice to say that there is a single well-designed interface
for class hierarchy queries in FindBugs.  However, this is not quite the
case.  The Subtypes class provides a nice interface for downwards queries.
The Hierarchy class provides a number of upwards query methods,
although these are somewhat less well thought-out.

\Note{
Designing and implementing a better API for class hierarchy queries 
would be a good idea.
}

\section{The Warning Object Model}

This section describes the object model used to represent FindBugs warnings.

\subsection{BugInstance}

When a detector finds code that is likely to be an error, it should
report a warning.  In FindBugs, the BugInstance class represents a single
warning, including the type of warning and associated context information.

\Note{
The BugInstance class really should have been called ``Warning''.
}

A BugInstance is created with three parameters:
\begin{description}
\item[Detector:] A reference to the Detector reporting the warning.
\item[Bug type:] This is a String that contains the warning's {\em bug code}.
\item[Priority:] One of HIGH\_PRIORITY, NORMAL\_PRIORITY,
                and LOW\_PRIORITY.  (These are integer constants that
                are defined in the Detector interface.)  By convention, normal or high
                priority is only used when there is a high degree of confidence
                that the user will consider the warning to represent a serious
                problem.
\end{description}

\Note{
NORMAL\_PRIORITY should have been called MEDIUM\_-PRIORITY.
}

\subsection{BugPatterns}

The bug type used to create a BugInstance
is a unique key that defines a particular {\em bug pattern}.
Behind the scenes, there is a BugPattern class that represents the
bug patterns defined by a particular FindBugs plugin.
Each bug pattern has some meta-information associated with it
that can be accessed through the BugPattern.

\subsection{BugAnnotation}

In its raw form, the BugInstance doesn't have any information about
important details such as what class the warning is being reported in,
what method, which source lines are affected, and so forth.
BugAnnotation objects provide this information.  Some kinds of BugAnnotations
include:
\begin{description}
\item[ClassAnnotation:] A class
\item[MethodAnnotation:] A method in a class
\item[FieldAnnotation:] A field of a class
\item[SourceLineAnnotation:] A range of source lines in a source file
\item[IntAnnotation:] An integer value
\end{description}

BugAnnotations should be added to a BugInstance in a well-defined order.
The order and types of BugAnnotations vary by bug pattern, but there are
some conventions.  The most important convention is that the first BugAnnotation
must be a ClassAnnotation: every warning is associated with a single class.
If the warning pertains to code, then the second BugAnnotation should be
a MethodAnnotation and the third should be a SourceLineAnnotation;
together, these define the code affected by the warning.  Otherwise,
BugAnnotations can be added in more or less free-form.  BugInstances sharing
a common bug type may even have different numbers of BugAnnotations,
depending on the exact circumstances in which the warnings were reported.

\subsection{I18N}

Formatting BugInstance and BugAnnotation objects in human readable
form is done using format strings.  The BugAnnotation format strings
are stored directly in findbugs.jar, and accessed through resource
bundles.  The BugInstance format strings are loaded through
XML files in the plugin containing the bug detector reporting the
warning.  The reason for this division is that BugAnnotations
are part of the core object model, and are not defined by detectors.
Bug patterns, on the other hand, are defined by individual plugins,
so their messages are local to the plugin.

A single BugAnnotation class can have multiple format strings,
depending on the {\em role} of the BugAnnotation.  For example,
a MethodAnnotation might be the method in which the warning is
reported (this is the default role), or could be a method called
at a call site reported by the method (the METHOD\_CALLED role).

Formatting model objects as human-readable Strings is done by
the I18N class.

\subsection{BugCollection}

The BugCollection class defines (no surprise) a collection of BugInstance
objects: generally, the output of a single analysis session.
In addition to storing BugInstance objects, the BugCollection also
stores some meta-information about the analysis session, such
as missing classes, reported analysis errors, etc.
In recent versions of FindBugs, a single BugCollection can represent
a history of analysis sessions on a single software artifact, where
each BugInstance may be present in multiple sessions.
The {\em workflow} tools can use this functionality to do data mining
on analysis results over time.

There is only one concrete BugCollection class, called
SortedBugCollection.

\Note{
BugCollection and SortedBugCollection are not a particularly elegant design.
}

\subsection{Reading and Writing XML}

The XML format supported by FindBugs is a very straightfoward representation of
the underlying object model.  As such, it is great for being able to store
a BugCollection and later load it without loss of information.  However,
because it does not contain human-readable messages, it is not very good
for conversion to other formats, e.g. via an XML stylesheet.

In recent versions, FindBugs supports an {\tt -xml:withMessages}
option that adds formatted strings (reflecting the current language/locale
settings) to the saved XML.  If the saved XML is re-loaded, these
messages will be discarded.  However, they are useful for creating
human-readable output (such as HTML reports) from the saved XML.

\end{document}

